% ---------------------------------------------------------------------------
% Theory subsection for the quality-gated farthest-point base selector
% (--sel_filter in final_update.py).
% Requires: amsmath, amssymb, amsthm, algorithm, algpseudocode.
% Theorem-like envs assumed:
%   \newtheorem{proposition}{Proposition}[section]
%   \newtheorem{remark}[proposition]{Remark}
% ---------------------------------------------------------------------------

\subsection{Quality-gated diverse base selection}
\label{sec:filter}

\paragraph{Setup and notation.}
Fix a target image $x_t$ with adversarial label $y_{\mathrm{adv}}$ and let
$\mathcal{C}=\{1,\dots,N\}$ index every training image of the poison class
($N=5000$ for CIFAR-10).  For a poison budget $\beta$ the attacker may perturb
$N_p=\lceil\beta n\rceil$ of them, so base selection is the combinatorial
problem of choosing $S\subseteq\mathcal{C}$ with $|S|=N_p$.  Let
$f_{\theta_1},\dots,f_{\theta_m}$ be the surrogate ensemble and
$\phi_j:\mathcal{X}\to\mathbb{R}^{d}$ the penultimate representation of
surrogate $j$.  Our selection score combines a target-proximity term
$d_j(i)$ and a confidence margin $\mu_j(i)$, each standardised over the
candidate pool by $z[\,\cdot\,]$ so that the two are commensurate:
\begin{equation}
  s(i)\;=\;\frac{1}{m}\sum_{j=1}^{m}\Bigl(z\bigl[d_j\bigr](i)
        \;+\;\lambda\, z\bigl[\mu_j\bigr](i)\Bigr),
  \qquad i\in\mathcal{C},
  \label{eq:score}
\end{equation}
with \emph{smaller $s$ meaning a more promising base}.  Alongside the score we
use the ensemble embedding obtained by $\ell_2$-normalising and concatenating
the per-surrogate features,
\begin{equation}
  \Psi(i)\;=\;\tfrac{1}{\sqrt{m}}\Bigl[\,
      \hat\phi_1(i)^{\!\top},\dots,\hat\phi_m(i)^{\!\top}\Bigr]^{\!\top}
      \in\mathbb{R}^{md},
  \qquad \hat\phi_j=\phi_j/\lVert\phi_j\rVert_2 .
  \label{eq:embed}
\end{equation}
The normalisation is chosen so that $\lVert\Psi(i)\rVert_2=1$ and
\begin{equation}
  \bigl\langle\Psi(i),\Psi(i')\bigr\rangle
  \;=\;\frac{1}{m}\sum_{j=1}^{m}\cos\!\bigl(\phi_j(i),\phi_j(i')\bigr),
  \label{eq:meancos}
\end{equation}
i.e.\ a \emph{single} inner product in the concatenated space equals the
\emph{mean} per-surrogate cosine similarity.  Two candidates are therefore
declared redundant only if every surrogate agrees that they are redundant,
which is what makes the criterion transfer to an unseen victim.  We write
$\rho(i,i')=1-\langle\Psi(i),\Psi(i')\rangle$ for the induced distance and,
for $S\subseteq\mathcal{A}\subseteq\mathcal{C}$,
\begin{equation}
  r(S;\mathcal{A})=\max_{i\in\mathcal{A}}\min_{j\in S}\rho(i,j),
  \qquad
  \operatorname{sep}(S)=\min_{i\neq j\in S}\rho(i,j)
  \label{eq:radii}
\end{equation}
for the covering radius of $S$ over $\mathcal{A}$ and the separation of $S$.

\paragraph{Why the score alone is not enough.}
The baseline selector takes the $N_p$ best-scoring candidates, i.e.\ it solves
\begin{equation}
  S^{\mathrm{plain}}\;=\;\operatorname*{arg\,min}_{|S|=N_p}\;F(S),
  \qquad F(S)=\sum_{i\in S}s(i).
  \label{eq:modular}
\end{equation}
$F$ is \emph{modular}: the marginal value $F(S\cup\{i\})-F(S)=s(i)$ of a
candidate never depends on what has already been chosen.  Modularity is
exactly the property that makes \eqref{eq:modular} trivially solvable by
sorting, and exactly the property that prevents the objective from expressing
``this candidate is already covered''.  Two consequences follow.

\emph{(i) Redundancy.}  Because $s$ is a smooth function of $\Psi$ through
\eqref{eq:score}, its sublevel sets are localised in embedding space.  The
top-$N_p$ set is the $N_p/N$ sublevel set of $s$, so it concentrates in
whichever lobe of the class manifold happens to lie closest to the target,
and near-duplicate images are selected together by construction.

\emph{(ii) Budget saturation.}  The ranking in \eqref{eq:modular} does not
depend on $N_p$.  As the budget grows, each additional slot is spent on the
next-best candidate by quality rather than on the candidate that contributes
the most \emph{new} information, so the selected set becomes an increasingly
biased sample of the class -- systematically low-margin, near-boundary images
-- while the quality signal that justified the bias becomes progressively less
discriminative.  This is the regime in which we observe score-based selection
underperforming a uniformly random base set.

\paragraph{Diversity as effective rank.}
The reason redundancy is harmful is visible in the crafting objective.  For
gradient matching the poison gradient of the selected set under perturbations
$\Delta=(\delta_i)_{i\in S}$ is, to first order,
\begin{equation}
  g_P(\Delta)\;\approx\;\bar{g}_S+\frac{1}{|S|}\sum_{i\in S}J_i\,\delta_i,
  \qquad
  J_i=\frac{\partial}{\partial x}\nabla_\theta
        L\bigl(f_\theta(x),y\bigr)\Big|_{x=x_i},
  \label{eq:linearise}
\end{equation}
so the attainable gradients form a zonotope inside the affine subspace
$\bar{g}_S+\operatorname{range}\bigl[J_1\;\cdots\;J_{N_p}\bigr]$.  Its
dimension -- the \emph{effective rank} of the base set -- upper-bounds how
closely $g_P$ can be steered towards the target gradient $g_t$: the best
achievable value of $1-\cos(g_P,g_t)$ is governed by the angle between $g_t$
and that subspace.  A duplicate base contributes $J_{i'}\approx J_i$ and hence
adds no new directions at all; under the $1/|S|$ normalisation it does not
even widen the zonotope.  The clean-label constraint
$\lVert\delta_i\rVert_\infty\le\varepsilon$ makes this binding, because the
attack must assemble a large component along $g_t$ out of many small,
\emph{distinct} contributions.  The same argument in feature space applies to
feature collision: co-located bases deform the victim's decision boundary in
one small region, and that deformation need not cover the augmentation-induced
spread of the target at test time.

This suggests replacing the modular objective by one that rewards coverage.
Two standard surrogates exist: the volumetric relaxation
$\log\det(I+\sigma^{-2}\Psi_S^{\top}\Psi_S)$, which is submodular and yields
the DPP-MAP selector, and the covering relaxation $r(S;\cdot)$, which yields a
$k$-centre problem.  We adopt the latter, constrained by the score.

\paragraph{The selector.}
Let $\gamma\ge1$ be a \emph{pool ratio} (\texttt{-{}-sel\_pool}, default
$\gamma=3$) and set
\begin{equation}
  k \;=\; \min\bigl(N,\;\max(N_p,\lceil\gamma N_p\rceil)\bigr),
  \qquad
  \mathcal{K} \;=\; \bigl\{\,i\in\mathcal{C}\;:\;s(i)\le s_{(k)}\,\bigr\},
  \label{eq:gate}
\end{equation}
where $s_{(k)}$ is the $k$-th order statistic of the scores.  The selector
then solves the $k$-centre problem \emph{restricted to the gate}
$\mathcal{K}$,
\begin{equation}
  S^{\mathrm{filter}}
  \;\approx\;\operatorname*{arg\,min}_{S\subseteq\mathcal{K},\,|S|=N_p}
      \; r(S;\mathcal{K}),
  \label{eq:kcenter}
\end{equation}
using farthest-point traversal seeded at the single best-scoring candidate.
Quality enters as a hard constraint rather than as a penalty term, and
diversity is optimised inside the feasible set.

\begin{algorithm}[t]
\caption{Quality-gated farthest-point base selection}
\label{alg:filter}
\begin{algorithmic}[1]
\Require scores $s\in\mathbb{R}^{N}$, embeddings $\Psi\in\mathbb{R}^{N\times md}$,
         budget $N_p$, pool ratio $\gamma\ge1$
\State $k\gets\min(N,\max(N_p,\lceil\gamma N_p\rceil))$
\State $\mathcal{K}\gets$ indices of the $k$ smallest entries of $s$
       \Comment{quality gate}
\State $j_1\gets\arg\min_{i\in\mathcal{K}}s(i)$;\quad
       $S\gets\{j_1\}$
\State $\mathrm{md}[i]\gets\rho(i,j_1)$ for all $i\in\mathcal{K}$
       \Comment{one mat-vec}
\For{$t=2,\dots,N_p$}
  \State $j_t\gets\arg\max_{i\in\mathcal{K}\setminus S}\mathrm{md}[i]$
  \State $S\gets S\cup\{j_t\}$
  \State $\mathrm{md}[i]\gets\min\bigl(\mathrm{md}[i],\rho(i,j_t)\bigr)$
         for all $i\in\mathcal{K}$
\EndFor
\State \Return $S$
\end{algorithmic}
\end{algorithm}

\paragraph{Properties.}
The construction has three properties that matter in practice.

\begin{proposition}[Exact interpolation]
\label{prop:interp}
If $\gamma=1$ then $S^{\mathrm{filter}}=S^{\mathrm{plain}}$.  If
$\gamma\ge N/N_p$ then the gate is vacuous and Algorithm~\ref{alg:filter}
reduces to unconstrained farthest-point traversal on $\mathcal{C}$.
\end{proposition}
\begin{proof}
For $\gamma=1$, \eqref{eq:gate} gives $k=N_p$, so $|\mathcal{K}|=N_p$ and the
loop must exhaust $\mathcal{K}$; the output equals the top-$N_p$ set
irrespective of the order in which points are added.  For $\gamma\ge N/N_p$,
$k=N$ and $\mathcal{K}=\mathcal{C}$.
\end{proof}

Proposition~\ref{prop:interp} is what makes the method a strict
generalisation rather than an alternative: the baseline is recovered exactly,
not approximately, at one end of the knob, and the pure-diversity selector at
the other.  Anything in between is a genuine trade-off, and $\gamma$ is
scale-free -- it is a ratio of set sizes, not a coefficient carrying the units
of the score.

\begin{proposition}[Bounded quality loss]
\label{prop:quality}
Every selected base lies in the top $k/N$ quantile of the score, and
\begin{equation}
  \frac{1}{N_p}\sum_{i\in S^{\mathrm{filter}}}\!\!s(i)
  \;-\;
  \frac{1}{N_p}\sum_{i\in S^{\mathrm{plain}}}\!\!s(i)
  \;\;\le\;\; s_{(k)}-s_{(N_p)} .
  \label{eq:qualbound}
\end{equation}
\end{proposition}
\begin{proof}
$S^{\mathrm{filter}}\subseteq\mathcal{K}$ gives
$\max_{i\in S^{\mathrm{filter}}}s(i)\le s_{(k)}$, and
$\min_{i\in S^{\mathrm{plain}}}s(i)=s_{(1)}\le s_{(N_p)}$ bounds each of the
$N_p$ paired differences by $s_{(k)}-s_{(N_p)}$.
\end{proof}

The bound in \eqref{eq:qualbound} is self-regulating.  The price paid for
diversity is a gap between two order statistics of the score distribution;
that gap is small precisely when the left tail of the score is flat, which is
also when the ranking carries the least information and diversity is most
worth buying.  Conversely, when the score is sharply peaked -- a few
genuinely privileged bases -- the gate is expensive to leave and the selector
stays close to the baseline.

\begin{proposition}[Coverage and anti-redundancy]
\label{prop:cover}
Let $r^\star=\min_{|S|=N_p,\,S\subseteq\mathcal{K}}r(S;\mathcal{K})$.
Algorithm~\ref{alg:filter} returns $S$ with
$r(S;\mathcal{K})\le 2r^\star$ and
$\operatorname{sep}(S)\ge r(S;\mathcal{K})$.
\end{proposition}
\begin{proof}
The $2$-approximation is Gonzalez's guarantee for farthest-point traversal
under a metric.  For the separation, let $r_t$ be the value of
$\max_i\mathrm{md}[i]$ at the start of iteration $t$; $\mathrm{md}$ is
pointwise non-increasing, so $r_2\ge r_3\ge\dots\ge r_{N_p}=r(S;\mathcal{K})$.
Each $j_t$ is at distance exactly $r_t\ge r_{N_p}$ from every previously
selected point, hence all pairwise distances in $S$ are at least
$r(S;\mathcal{K})$.
\end{proof}

The separation bound is the formal statement of the anti-redundancy property:
no two selected bases can be near-duplicates, and the guarantee degrades
gracefully as the budget grows (larger $N_p$ forces a smaller
$r(S;\mathcal{K})$, hence a smaller admissible separation).

\begin{remark}[Budget adaptivity]
\label{rem:adaptive}
The clipping in \eqref{eq:gate} makes the gate automatically budget-dependent:
it is active iff $\gamma N_p<N$, i.e.\ iff
$\beta<\beta^\star=N/(\gamma n)$.  With $N=5000$, $n=50\,000$ and $\gamma=3$,
$\beta^\star\approx3.3\times10^{-2}$: at $\beta=5\times10^{-3}$ the gate
retains $k=750$ of $5000$ candidates and the selector is a true
quality--diversity blend, whereas at $\beta=4\times10^{-2}$ the gate is
vacuous and the selector degenerates to pure coverage.  The score is thus
weighted most heavily where it is most reliable -- small $N_p$, where the
top of the ranking is sharply better than the bulk -- and is annealed away as
the budget approaches the class size, without any budget-specific tuning.
\end{remark}

\begin{remark}[Cost]
\label{rem:cost}
Stage one is a partial sort, $O(N\log k)$.  Stage two maintains the
running minimum $\mathrm{md}$, so each of the $N_p$ iterations costs one
mat-vec against the gated embeddings: $O(N_p k m d)$ time and $O(k m d)$
memory.  Both scale linearly in $k$, in contrast to the $O(N^2)$ kernel
required by a marginal-relevance or determinantal selector over the full
candidate pool, and the gate shrinks $k$ exactly in the small-budget regime
where the largest number of greedy iterations would otherwise be cheapest to
avoid.  Selection is performed under \texttt{torch.no\_grad()} and adds a
negligible fraction of end-to-end crafting time.
\end{remark}

\paragraph{Relation to scalarised alternatives.}
Marginal-relevance and determinantal selectors optimise a single scalarised
objective, $s(i)+\mu\max_{j\in S}\langle\Psi(i),\Psi(j)\rangle$ and
$\log\det L_S$ with quality $q_i=\exp(-\alpha s(i))$ respectively.  Both are
principled, but their trade-off coefficients are not scale-free: $\mu$ carries
the units of the score, and $\alpha$ interacts with the spectrum of the
similarity kernel, so a value tuned at one budget or one architecture need not
transfer.  (We mitigate this for the marginal-relevance variant by
standardising the similarity term against the empirical distribution of
pairwise similarities, and for the determinantal variant by shifting the
quality by $\min_i s(i)$ before exponentiating, which is required for
numerical stability in the large-$\alpha$ limit where it must recover the
baseline.)  The constrained formulation \eqref{eq:kcenter} avoids the issue
altogether: it has one dimensionless knob, an exact reduction to the baseline
at $\gamma=1$, and the budget-adaptive behaviour of
Remark~\ref{rem:adaptive} comes for free.
